\documentclass[../../main.tex]{subfiles}
\begin{document}

$A$を2次の正方行列とする．
列ベクトル$\vec{x}_0$に対し，列ベクトル$\vec{x}_1$，$\vec{x}_2$，$\cdots$を
$\vec{x}_{n+1}=A\vec{x}_n$ $(n=0,1,2,\cdots)$によって定める．
ある零ベクトルではない$\vec{x}_0$について，3以上の自然数$m$で初めて$\vec{x}_m$が$\vec{x}_0$と一致するとき，
行列$A^m$は単位行列であることを示せ．

\end{document}
